\documentclass[11pt,letterpaper,sans]{moderncv}

\moderncvstyle{classic}
\moderncvcolor{purple}
\moderncvicons{marvosym}
% Keep the name uniformly black across moderncv versions, including Overleaf.
\colorlet{namecolor}{black}
\colorlet{firstnamecolor}{black}
\colorlet{lastnamecolor}{black}

\usepackage[utf8]{inputenc}
\usepackage[T1]{fontenc}
\usepackage{lmodern}
\usepackage[top=0.5in,bottom=0.5in,left=0.5in,right=0.5in]{geometry}
\setlength{\hintscolumnwidth}{3.3cm}
\urlstyle{same}
\nopagenumbers{}
% Keep terms and DOIs intact for applicant-tracking text extraction
\hyphenpenalty=10000
\exhyphenpenalty=10000
\tolerance=2000
\emergencystretch=3em

\name{Bohan}{Lu}
\title{PhD Candidate, Quantum Error Correction}
\email{bohan.lu@duke.edu}
\homepage{jerrylvx.github.io}
\extrainfo{\httpslink{github.com/jerrylvx}}

\begin{document}

% Plain-text contact line and bullets (icon glyphs extract as stray letters)
\renewcommand*{\emailsymbol}{}
\renewcommand*{\homepagesymbol}{}
\renewcommand*{\listitemsymbol}{\textbullet~}
\renewcommand{\labelitemi}{\textcolor{color1}{\textbullet}}
% Name at 24 pt (Duke guide: 16 pt or larger)
\renewcommand*{\namefont}{\fontsize{24}{26}\mdseries\upshape}

\makecvtitle
\vspace{-3em}

\section{Research Interests}
\cvitem{}{Fault-tolerant quantum computing, especially transversal non-Clifford gates, quantum LDPC and Floquet codes, and decoding. Recent work uses projective geometry and SAT to construct codes implementing logical CCZ through transversal physical $T/T^\dagger$ gates.}

\section{Education}
\cventry{2023--present}{PhD in Electrical and Computer Engineering}{Duke University}{Durham, NC}{expected 2027--2028}{\normalsize Advisor: Kenneth R. Brown}
\cventry{2019--2022}{MSc in Physics}{Freie Universität Berlin}{Berlin, Germany}{}{\normalsize Thesis: \emph{Entanglement dynamics of a many-body localized system}. Advisor: Jens Eisert}
\cventry{2015--2019}{BA in Mathematics and Physics}{Washington University in St.\ Louis}{St.\ Louis, MO}{}{\normalsize Cum Laude, Distinction in Mathematics, High Distinction in Physics}

\section{Research Experience}
\cventry{2023--present}{Doctoral Research, Brown Group}{Duke University}{Durham, NC}{}{\normalsize
\begin{itemize}
\item Constructed CSS codes implementing logical CCZ via physical transversal $T/T^\dagger$ using projective geometry and SAT, and proved $n\geq 39$ for CSS codes with 3 logical qubits, $d_Z\geq 3$, and quasi-transversal CCZ
\item Designed a dynamical code from a minimally non-local hypergraph product code
\item Developed stim-floquet (\href{https://github.com/jerrylvx/stim-floquet}{github.com/jerrylvx/stim-floquet}), an open-source Python package for Floquet schedules, GF(2) homology checks, local detector discovery, and Stim circuit output
\end{itemize}}
\cventry{\parbox[t]{\hintscolumnwidth}{\raggedleft Oct 2025--May 2026\\Aug 2026--present}}{Consulting Contractor}{Xanadu Quantum Technologies}{Remote}{}{\normalsize
\begin{itemize}
\item Research fault-tolerant quantum architecture (supervisor: Eric Sabo)
\end{itemize}}
\cventry{May--Aug 2026}{Architecture Intern}{Xanadu Quantum Technologies}{Remote}{}{\normalsize
\begin{itemize}
\item Researched decoding for quantum LDPC codes (supervisor: Eric Sabo)
\end{itemize}}
\cventry{Jul--Oct 2022}{Research Intern}{HQS Quantum Simulations GmbH}{Karlsruhe, Germany}{}{\normalsize
\begin{itemize}
\item Developed and benchmarked a Hermite-polynomial spectral method for resonant-level models and wrote an internal report (supervisor: Peter Schmitteckert)
\end{itemize}}
\cventry{2021--2022}{Master's Thesis Research}{Freie Universität Berlin}{Berlin, Germany}{}{\normalsize
\begin{itemize}
\item Co-developed an entanglement-based method to extract spatial properties of quasi-local integrals of motion, published in Communications Physics
\item Co-wrote simulation code and ran tensor-network simulations on clusters
\end{itemize}}

\section{Publications}
\cvitem{2026}{\begin{minipage}[t]{\linewidth}\raggedright\textbf{Building codes with transversal CCZ using projective geometry and SAT solvers.} \textbf{B.~Lu}, K.~R.~Brown. Preprint, \href{https://arxiv.org/abs/2610.09341}{arXiv:2610.09341}. Code and data: \href{https://github.com/jerrylvx/gate2code-public}{github.com/jerrylvx/gate2code-public}\end{minipage}}
\cvitem{2024}{\begin{minipage}[t]{\linewidth}\raggedright\textbf{Measuring out quasi-local integrals of motion from entanglement.} \textbf{B.~Lu}, C.~Bertoni, S.~J.~Thomson, J.~Eisert. \emph{Communications Physics} \textbf{7}, 17 (2024). \href{https://doi.org/10.1038/s42005-023-01478-5}{doi:10.1038/s42005-023-01478-5}, \href{https://arxiv.org/abs/2301.01787}{arXiv:2301.01787}\end{minipage}}

\section{Skills}
\cvitem{AI skills}{Agentic orchestration, engineering, and harness building}
\cvitem{Computing}{Python, C++, HPC clusters}

\clearpage
\section{Talks}
\cvitem{Oct 2026}{\emph{Building codes with transversal CCZ using projective geometry and SAT solvers.} \href{https://seminar.math.vt.edu/applied_algebra/index.html}{Applied Algebra Seminar}, Virginia Tech (online), Oct 22, 2026 (upcoming).}
\cvitem{Sep 2026}{\emph{Boolean functions: applications in quantum error correction and beyond.} \href{https://sites.duke.edu/dqctheory/}{Duke Quantum Center Theory Seminar}, Sep 16, 2026.}
\cvitem{2019}{\emph{Reduced density operators of finite-sized Heisenberg XXZ chains.} With A.~Seidel. APS March Meeting 2019.}

\section{Posters}
\cvitem{2026}{\emph{Discovery and classification of small CSS codes supporting a quasi-transversal CCZ logical gate.} QEC 2026.}
\cvitem{2025}{\emph{A dynamical error-correcting code from a minimally non-local hypergraph product code.} With K.~Brown. QIP 2025 and QEC 2025.}
\cvitem{2023}{\emph{Improved decoding of quantum LDPC codes using neural-network enhanced belief propagation.} With A.~Pesah, J.~Roffe, N.~Raveendran. QIP 2023, QEC 2023 (Sydney), and \emph{Coping with Errors in Scalable Quantum Computing Systems}, WE-Heraeus Seminar (Bad Honnef).}
\cvitem{2022}{\emph{Entanglement dynamics in a many-body localized system.} With C.~Bertoni, S.~Thomson, J.~Eisert. CRC 183 Summer School 2022.}

\section{Teaching}
\cventry{Spring 2026}{Teaching Assistant, Liberata}{Duke Bass Connections}{Duke University}{}{\normalsize
\begin{itemize}
\item Mentored student research projects and contributed to research and development on the algorithms team
\end{itemize}}
\cvitem{Fall 2024}{Teaching Assistant, ECE 520 Quantum Engineering, Duke University}
\cvitem{Winter 2020}{Teaching Assistant, tensor-network seminar, Freie Universität Berlin}

\section{Awards and Certificates}
\cvitem{2026}{Duke Graduate School Conference Travel Award}
\cvitem{2025}{NVIDIA Academy certificate, AI Infrastructure and Operations Fundamentals}
\cvitem{2017--2018}{Greg Delos Fellowship, Physics Department, Washington University in St.\ Louis}
\cvitem{2016}{Undergraduate Research Fellowship, Physics Department, Washington University in St.\ Louis}
\cvitem{2015--2019}{Dean's List, Washington University in St.\ Louis (Fall 2015 to Fall 2016, Spring 2019)}

\end{document}
